%% empirical_design_draft.tex
%% Standalone draft — not yet integrated into main paper.
%% Flag to Neil: variable definitions should be cross-checked against
%% final constructed values in scripts 01–04 before merging.

\section{Empirical Design}
\label{sec:design}

\subsection{Sample Construction}

\textbf{Auditor change events.} We identify auditor changes using Form 8-K
filings disclosing a change in the registrant's certifying accountant (Item
4.01). Item 4.01 disclosures are mandatory within four business days of the
change and must include the name of the departing accountant, whether the
departure was a dismissal or resignation, and whether any disagreements arose
over accounting or disclosure matters. We obtain these filings from SEC EDGAR
for the period 2000 to 2023, yielding an initial sample of approximately
[PLACEHOLDER: N] events.

We retain only original 8-K filings (not 8-K/A amendments) and require the
filing to be matched to a CRSP PERMNO via the CRSP-Compustat linking table.
We further require at least 100 non-missing trading days in the estimation
window (see below). The final sample consists of [PLACEHOLDER: N] auditor
change events for [PLACEHOLDER: N] unique firms.

\textbf{Exclusions.} Following standard practice in the auditor change
literature \citep{defond2014}, we exclude financial firms (SIC 6000--6999)
and utilities (SIC 4900--4999), whose audit relationships are subject to
additional regulatory constraints. We also exclude the first year of EDGAR
electronic filing coverage for each firm to avoid systematic filing quality
differences in early years.

\subsection{Dependent Variables}

\textbf{Absolute cumulative abnormal return ($|CAR|$).} We estimate a
market model for each firm using daily returns over the estimation window
$[-252, -46]$ trading days relative to the event date, regressing firm
returns on the CRSP value-weighted market return. The market model
parameters $\hat{\alpha}$ and $\hat{\beta}$ are then used to compute
abnormal returns over the $[-1, +1]$ event window:
\begin{equation}
    AR_{i,t} = R_{i,t} - (\hat{\alpha}_i + \hat{\beta}_i \cdot R_{m,t}),
    \quad t \in \{-1, 0, +1\}.
\end{equation}
The cumulative abnormal return is $CAR_i = \sum_{t=-1}^{+1} AR_{i,t}$, and
the primary dependent variable is $|CAR_i|$. We include delisting returns
following \citet{beaver2007} to mitigate survivorship bias.

\textbf{Abnormal trading volume ($AbVol$).} For each firm, we estimate mean
log-volume over the estimation window $[-252, -46]$ relative to the event
date. Abnormal volume in the event window is the mean of
$[\log(\text{Vol}_{i,t} + 1) - \overline{\log(\text{Vol}_i)}]$ over
$t \in \{-1, 0, +1\}$.

\subsection{Political Polarization Measure}

We measure political polarization using the \citet{esteban1994} index applied
to U.S. House of Representatives election returns from the MIT Election Lab.
For congressional district $d$ in election year $e$, let $s_{D,d,e}$ and
$s_{R,d,e}$ denote the Democratic and Republican shares of the two-party vote.
The district-level polarization index is:
\begin{equation}
    \label{eq:er}
    P^{ER}_{d,e} = s_{D,d,e} \cdot s_{R,d,e}
    \cdot \left(s_{D,d,e}^\alpha + s_{R,d,e}^\alpha\right),
    \quad \alpha = 1,
\end{equation}
which simplifies to $s_{D,d,e} \cdot s_{R,d,e}$ when $\alpha = 1$ (since
$s_D + s_R = 1$ on the two-party vote share). This measure is maximized at
0.25 when the two parties split the vote equally, and decreases as the vote
becomes more one-sided. We aggregate to the state level by weighting each
district by its share of total state two-party votes:
\begin{equation}
    P^{ER}_{s,e} = \sum_{d \in s}
    \frac{\text{TwoPartyVotes}_{d,e}}{\text{TwoPartyVotes}_{s,e}}
    \cdot P^{ER}_{d,e}.
\end{equation}
Since House elections occur biennially, we forward-fill each election-year
measure to the subsequent odd year, giving an annual state-by-year panel.
We match firms to their headquarters state using the Compustat \texttt{state}
variable and merge on filing year. We standardize $P^{ER}_{s,t}$ to have
mean zero and unit variance to facilitate coefficient interpretation.

The ER measure has a natural behavioral microfoundation: it captures the
degree to which the population is divided into internally homogeneous,
mutually alienated camps \citep{esteban1994}. States closer to a 50-50
partisan split have higher ER scores, reflecting greater polarization between
two equally sized ideological groups — precisely the configuration expected to
generate the most disagreement in posterior beliefs.

\subsection{Empirical Specification}

Our main specification is:
\begin{equation}
    \label{eq:main}
    Y_{i,t} = \alpha + \beta_1\,\textit{Event}_{i,t}
    + \beta_2\,\textit{Pol}_{s(i),t}
    + \beta_3\,(\textit{Event}_{i,t} \times \textit{Pol}_{s(i),t})
    + \gamma' X_{i,t} + \delta_i + \tau_t + \varepsilon_{i,t},
\end{equation}
where $Y_{i,t}$ is either $|CAR_{i,t}|$ or $AbVol_{i,t}$;
$\textit{Event}_{i,t}$ is an indicator equal to one in the event window;
$\textit{Pol}_{s(i),t}$ is the standardized ER polarization measure for firm
$i$'s headquarters state $s(i)$ in year $t$; $X_{i,t}$ is a vector of
firm-level controls (size, leverage, ROA, book-to-market, loss indicator,
sales growth); and $\delta_i$ and $\tau_t$ are firm and year fixed effects.
Standard errors are clustered at the firm level.

The coefficient of interest is $\beta_3$, the interaction of the event
indicator with polarization. A positive and significant $\beta_3$ indicates
that polarization amplifies the market reaction to auditor change disclosures,
consistent with Hypothesis~\ref{h1}.

\subsection{Identification}

The key identification challenge is that polarization may be correlated with
other state-level economic or institutional conditions that independently
affect market reactions. We address this in several ways.

First, firm fixed effects absorb all time-invariant firm-level confounders,
including stable characteristics such as industry membership, auditor
relationships, and listing exchange. Year fixed effects control for aggregate
market conditions that affect all firms in a given year.

Second, the interaction design provides an additional layer of identification:
$\beta_3$ identifies the \textit{differential} effect of polarization on
event-window returns relative to non-event windows, within the same firm
over time. A confound would need to predict not only higher returns in
polarized states but specifically higher returns \textit{around auditor change
events} in polarized states.

Third, we exploit time-series variation in polarization within states. The ER
measure changes from election cycle to election cycle as vote shares shift. A
firm headquartered in a state that becomes more polarized over time experiences
stronger market reactions to subsequent auditor changes. This within-state
variation is unlikely to be driven by firm-level reporting decisions.

\subsection{Controls}

We include the following firm-year controls in $X_{i,t}$:
\begin{itemize}
    \item \textit{Size}: natural log of total assets; larger firms have more
          analyst coverage and potentially less information asymmetry.
    \item \textit{Leverage}: total debt divided by total assets; more
          levered firms face greater financial distress risk, which may
          independently affect reactions to audit signals.
    \item \textit{ROA}: net income divided by total assets; controls for
          underlying firm profitability.
    \item \textit{Book-to-market}: book equity divided by market equity;
          controls for value vs.\ growth characteristics.
    \item \textit{Loss}: indicator equal to one if net income is negative;
          loss firms are more likely to switch auditors for distress-related
          reasons.
    \item \textit{Sales growth}: year-over-year revenue growth; controls
          for firm lifecycle and organic change in audit complexity.
\end{itemize}
All continuous controls are winsorized at the 1st and 99th percentiles.
